\vfill\pagebreak
\section{Declaring and calling a function}
\label{sec:functions:declaring}

The report below prints a line of dashes three times. Rather than writing the
same \token{println} three times, it declares a function \token{separator} that
prints the line, and calls it wherever a line is needed.

\lstinputlisting[style=coloredverbatim, caption={A function called three times, \textit{examples/functions/separator.yr}}, label=lst:functions:separator]{examples/functions/separator.yr}
\vspace{-10pt}%
\begin{lstlisting}[style=bashVerb, escapechar=@+]
@+\textcolor{teal!80}{alice@dev:\mtilde}@+$ gyc separator.yr -o separator
@+\textcolor{teal!80}{alice@dev:\mtilde}@+$ ./separator
--------------------
  Monthly report
--------------------
Sales: 1250
Costs:  830
--------------------
\end{lstlisting}

A function declaration has four parts, which lines~3 to~5 show in order:

\begin{itemize}
\item the keyword \token{fn}, which announces a function;
\item its name, \token{separator}, chosen by the programmer under the rules of
  identifiers (Section~\ref{sec:types:constants_variables}). By convention the name
  starts with a lowercase letter, and a name of several words capitalizes each
  word after the first, as in \token{printCalendar};
\item a pair of parentheses, which holds the function's parameters -- here there
  are none, the next subsection gives it some;
\item its \textit{body}, a block (Section~\ref{sec:control:blocks}) that holds
  the instructions to run.
\end{itemize}

A declaration runs nothing by itself. The body runs when the function is
\textit{called}, by writing its name followed by parentheses, as on lines~8,
10 and 13. A call jumps to the start of the body, runs it to the end, and then
comes back to the caller, just after the call, where the program carries on.
Figure~\ref{fig:functions:call_return} follows the cursor through the first two
calls. \token{main} itself is a function, the one the program calls first, and
\token{println} is a function too, declared in the module \token{std::io}.

\begin{figure}[h]
  \centering
  \begin{adjustbox}{max width=\linewidth}
    \begin{tikzpicture}[node distance=6mm]

      \node[cfterm] (start) {start of main};
      \node[cfstep, below=of start] (call1) {separator();};
      \node[cfstep, below=12mm of call1] (title) {println("  Monthly report");};
      \node[cfstep, below=of title] (call2) {separator();};
      \node[cfstep, below=12mm of call2] (sales) {println("Sales: 1250");};
      \node[cflab, below=4mm of sales] (more) {\dots};

      \node[cfstep, right=22mm of call1, yshift=-9mm] (body1) {println("-----...");};
      \node[cfstep, right=22mm of call2, yshift=-9mm] (body2) {println("-----...");};
      \node[cflab, right=2mm of body1, align=left] {body of\\separator};
      \node[cflab, right=2mm of body2, align=left] {body of\\separator};

      \draw[cfflow] (start) -- (call1);
      \draw[cfflow] (call1.east) -| node[cflab, above, pos=0.25] {call} (body1.north);
      \draw[cfflow] (body1.south) |- node[cflab, below, pos=0.75] {return} (title.east);
      \draw[cfflow] (title) -- (call2);
      \draw[cfflow] (call2.east) -| node[cflab, above, pos=0.25] {call} (body2.north);
      \draw[cfflow] (body2.south) |- node[cflab, below, pos=0.75] {return} (sales.east);
      \draw[cfflow] (sales) -- (more);
    \end{tikzpicture}
  \end{adjustbox}
  \caption{\label{fig:functions:call_return} The first two calls of
    Listing~\ref{lst:functions:separator}. Each call runs the body of
    \token{separator} from the start, then control returns to the statement that
    follows the call.}
\end{figure}


If the separator should become a line of stars, one line of the program changes,
line~4, and all three separators follow. With three copies of the
\token{println}, all three would have to be found and changed, and missing one
would leave a report whose lines do not match.

As Chapter~\ref{chap:basics} said, the order in which functions are declared in
the file does not matter: \token{main} may call a function declared below it.

\subsection{Parameters}
\label{sec:functions:parameters}

A function that always does exactly the same thing is of limited use. A function
drawing a line of stars is more useful if each call can say how long the line
is. A \textit{parameter} is a variable of the function whose value is given by
each call. It is declared between the parentheses of the declaration, as a name
followed by a colon and a type, like a variable declaration without its
\token{let}. The value that a call gives it, written between the parentheses of
the call, is called an \textit{argument}.

\lstinputlisting[style=coloredverbatim, caption={Functions with parameters, \textit{examples/functions/rectangle.yr}}, label=lst:functions:rectangle]{examples/functions/rectangle.yr}
\vspace{-10pt}%
\begin{lstlisting}[style=bashVerb]
************

******
******
******
\end{lstlisting}

\token{line} has one parameter, \token{width}, of type \token{i32}. The call
\token{line(12)} on line~17 runs the body of \token{line} with \token{width}
holding 12, and each call gets the value it passes: \token{width} is 6 during the
calls that \token{rectangle} makes.

\token{rectangle} has two parameters, separated by a comma, and a call gives
them two arguments in the same order: in \token{rectangle(6, 3)}, \token{width}
is 6 and \token{height} is 3. It also shows that a function can call another
function. \token{rectangle} does not know how a line is drawn, only that
\token{line} draws one, so a rectangle is three lines of 6 stars.

An argument can be any expression of the right type, not only a literal:
\token{line(n)}, \token{line(n * 2 + 1)} or \token{line(read!i32())} are all
calls to \token{line}. The expression is computed first, then its value is given
to the parameter.

\subsection{Parameters are constants}

A parameter holds a copy of its argument, and the function cannot change it:
like a variable declared without \token{mut}
(Section~\ref{sec:types:variable_declaration}), a parameter is a constant. The
next listing tries to count down with its parameter, and is refused.

\begin{lstlisting}[style=coloredverbatimError, escapechar=@, caption=Assigning a parameter]
use std::io;

fn countdown(from: i32) {
  while from > 0 {
    println(from);
    @\hb{from}@ -= 1; // error
  }
}

fn main() {
  countdown(3);
}
\end{lstlisting}

The compiler says \errmsg{E4156}{from is a value parameter, and cannot be used as
  a lvalue},
an \textit{lvalue} being something that can stand on the left of an assignment.
A function that needs a value it can change declares a variable of its own,
initialized from the parameter.

\begin{lstlisting}[style=coloredverbatim, escapechar=@, caption=A mutable copy of a parameter]
use std::io;

fn countdown(from: i32) {
  @\hcb{let mut n = from;}@
  while n > 0 {
    println(n);
    n -= 1;
  }
}

fn main() {
  countdown(3);
}
\end{lstlisting}

A parameter follows the other rules of variables, too. One that the body never
reads is refused, as an unused variable is
(Section~\ref{sec:types:variable_declaration}). A parameter that must exist
without being read, because every call passes an argument for it, is named
\token{\_}, and the value passed for it is ignored. A name that starts and ends
with an underscore, such as \token{\_width\_}, is not checked either, but it
still says what the argument is, and the body can still read it. Both forms
work for any variable (Section~\ref{sec:memory:unused_variables}).

\subsection{Matching the arguments}

A call must give exactly one argument per parameter, and each argument must have
the type of its parameter. The call converts nothing. A literal such as
\token{12} takes the type the parameter asks for, but a variable must already
have that type.

\begin{lstlisting}[style=coloredverbatimError, escapechar=@, caption=Arguments that do not match the parameters]
use std::io;

fn line(width: i32) {
  for _ in 0 .. width {
    print("*");
  }
  println();
}

fn main() {
  line(12);    // ok, the literal is an i32
  let n = 12u8;
  line(@\hb{n}@);     // error, n is a u8
  line(@\hb{3, 4}@);  // error, one argument too many
}
\end{lstlisting}

Both calls are refused with the same message, \errmsg{E4226}{the call operator
  is not defined for \errhole{} and \errhole{}}, the holes being the function
and the arguments of the call. Under it, the compiler lists each function that
the call could have meant, and why it does not fit:
\errmsg{E4095}{incompatible types i32 and u8} for the first,
\errmsg{E4212}{1 parameters is expected, but 2 are provided} for the second.
